% !TEX root = ../LosningsforslagTRES0312.tex
\chapter{Arbeid og energi}
\label{LF:Chapter:energi}
Dette kapittelet inneholder løsningsforslag til oppgavene i Kapittel~\ref{komp-Chapter:energi}, \textit{Arbeid og energi}, i kompendiet.

\oppgaver{Arbeid med konstant kraft}{k5oppg1}{%
%
Eva drar en kasse $s = \SI{12}{\metre}$ bortover et flatt gulv. Hun drar i et tau med kraften $F = \SI{150}{\newton}$, og tauet danner vinkelen $\theta = 25^\circ$ med gulvet. Friksjonskraften på kassen er $R = \SI{110}{\newton}$.
\begin{enumerate}[a)]
    \item Hvor stort arbeid gjør Eva på kassen?
    \item Hvor stort arbeid gjør friksjonskraften? Forklar hvorfor svaret blir negativt.
    \item Hvor stort arbeid gjør tyngdekraften og normalkraften på kassen? Forklar svaret.
\end{enumerate}
%
}

\svar[title={Løsningsforslag}]{%
\noindent
Vi bruker definisjonen av arbeid, $W = F\cdot s\cdot\cos\theta$, der $\theta$ er vinkelen mellom kraften og forflytningen (se \textit{Arbeid} i Kapittel~\ref{komp-Chapter:energi}).

\smallskip\noindent
\textbf{a)} Tauet danner vinkelen $\theta = 25^\circ$ med bevegelsesretningen:
\[
    W_F = F\cdot s\cdot\cos 25^\circ = \SI{150}{\newton}\cdot\SI{12}{\metre}\cdot 0.906 = \SI{1.63}{\kilo\joule} \approx \doubleunderline{\SI{1.6}{\kilo\joule}}
\]

\smallskip\noindent
\textbf{b)} Friksjonskraften peker motsatt vei av bevegelsen, så $\theta = 180^\circ$ og $\cos\theta = -1$:
\[
    W_R = R\cdot s\cdot\cos 180^\circ = \SI{110}{\newton}\cdot\SI{12}{\metre}\cdot(-1) = \SI{-1.32}{\kilo\joule} \approx \doubleunderline{\SI{-1.3}{\kilo\joule}}
\]
Arbeidet er negativt fordi friksjonen virker \emph{mot} bevegelsen. Friksjonen tar energi ut av bevegelsen og gjør den om til varme.

\smallskip\noindent
\textbf{c)} Tyngdekraften peker rett ned og normalkraften rett opp. Begge står vinkelrett på forflytningen ($\theta = 90^\circ$, $\cos 90^\circ = 0$), så
\[
    \doubleunderline{W_G = 0} \qquad\text{og}\qquad \doubleunderline{W_N = 0}
\]
Kassen flytter seg verken opp eller ned, så disse kreftene bidrar ikke til bevegelsen.
}

\oppgaver{Arbeid på en fjær}{k5oppg2}{%
%
En fjær har fjærkonstant $k = \SI{400}{\newton\per\metre}$.
\begin{enumerate}[a)]
    \item Hvor stort arbeid må du gjøre for å presse fjæra sammen $x = \SI{5.0}{\centi\metre}$?
    \item Hvor stort arbeid må du gjøre for å presse den sammen $x = \SI{10}{\centi\metre}$?
    \item Kompresjonen ble doblet, men arbeidet ble ikke doblet. Forklar hvorfor, gjerne ved hjelp av arealet under $F(x)$-grafen.
\end{enumerate}
%
}

\svar[title={Løsningsforslag}]{%
\noindent
Arbeidet for å presse sammen en fjær er arealet under $F(x)$-grafen, $W = \tfrac{1}{2}kx^2$.

\smallskip\noindent
\textbf{a)} Med $x = \SI{0.050}{\metre}$:
\[
    W = \frac{1}{2}\cdot\SI{400}{\newton\per\metre}\cdot(\SI{0.050}{\metre})^2 = \doubleunderline{\SI{0.50}{\joule}}
\]

\smallskip\noindent
\textbf{b)} Med $x = \SI{0.10}{\metre}$:
\[
    W = \frac{1}{2}\cdot\SI{400}{\newton\per\metre}\cdot(\SI{0.10}{\metre})^2 = \doubleunderline{\SI{2.0}{\joule}}
\]

\smallskip\noindent
\textbf{c)} Arbeidet ble fire ganger så stort. Arbeidet avhenger av $x^2$: dobles $x$, blir $x^2$ fire ganger så stort. Grafisk er arbeidet arealet av en trekant under $F(x) = kx$. Når $x$ dobles, blir både grunnlinjen og høyden i trekanten dobbelt så store, og da blir arealet fire ganger så stort.
}

\oppgaver{Kinetisk energi}{k5oppg3}{%
%
\begin{enumerate}[a)]
    \item En bil med masse $m = \SI{1200}{\kilo\gram}$ kjører i \SI{50}{\kilo\metre\per\hour}. Finn den kinetiske energien til bilen.
    \item Finn den kinetiske energien når bilen kjører i \SI{100}{\kilo\metre\per\hour}. Hvor mange ganger større er den? Forklar ved hjelp av formelen for $E_k$.
    \item En sprinter med masse \SI{70}{\kilo\gram} løper med fart \SI{10}{\metre\per\second}, og en geværkule med masse \SI{10}{\gram} har fart \SI{900}{\metre\per\second}. Hvem har mest kinetisk energi?
\end{enumerate}
%
}

\svar[title={Løsningsforslag}]{%
\noindent
Vi bruker $E_\text{K} = \tfrac{1}{2}mv^2$, og gjør om farten til \si{\metre\per\second} ved å dele på 3.6.

\smallskip\noindent
\textbf{a)} $v = \SI{50}{\kilo\metre\per\hour} = \SI{13.9}{\metre\per\second}$:
\[
    E_\text{K} = \frac{1}{2}\cdot\SI{1200}{\kilo\gram}\cdot(\SI{13.9}{\metre\per\second})^2 = \SI{116}{\kilo\joule} \approx \doubleunderline{\SI{1.2e2}{\kilo\joule}}
\]

\smallskip\noindent
\textbf{b)} $v = \SI{100}{\kilo\metre\per\hour} = \SI{27.8}{\metre\per\second}$:
\[
    E_\text{K} = \frac{1}{2}\cdot\SI{1200}{\kilo\gram}\cdot(\SI{27.8}{\metre\per\second})^2 = \doubleunderline{\SI{463}{\kilo\joule}}
\]
Det er \doubleunderline{fire ganger} så mye. Siden $E_\text{K}$ er proporsjonal med $v^2$, gir dobbel fart $2^2 = 4$ ganger så mye kinetisk energi. Det er derfor en kollisjon i \SI{100}{\kilo\metre\per\hour} er så mye farligere enn i \SI{50}{\kilo\metre\per\hour}.

\smallskip\noindent
\textbf{c)}
\[
    E_\text{K,sprinter} = \frac{1}{2}\cdot\SI{70}{\kilo\gram}\cdot(\SI{10}{\metre\per\second})^2 = \SI{3.5}{\kilo\joule}
\]
\[
    E_\text{K,kule} = \frac{1}{2}\cdot\SI{0.010}{\kilo\gram}\cdot(\SI{900}{\metre\per\second})^2 = \SI{4.05}{\kilo\joule}
\]
\doubleunderline{Geværkula} har litt mer kinetisk energi enn sprinteren, selv om massen er 7000 ganger mindre. Den store farten veier opp for den lille massen.
}

\oppgaver{Kasse som sklir til ro}{k5oppg4}{%
%
En kasse med masse $m = \SI{20}{\kilo\gram}$ sklir bortover et gulv med startfart $v_0 = \SI{3.0}{\metre\per\second}$. Den eneste kraften som bremser kassen, er friksjonskraften $R = \SI{40}{\newton}$.
\begin{enumerate}[a)]
    \item Hvor stor er endringen i kinetisk energi fra kassen begynner å skli til den står stille?
    \item Bruk arbeid-energi-setningen til å finne hvor langt kassen sklir før den stopper.
    \item Hvor langt ville kassen skli dersom startfarten var dobbelt så stor?
\end{enumerate}
%
}

\svar[title={Løsningsforslag}]{%
\noindent
\textbf{a)} Kassen stopper, så den kinetiske energien går fra $\tfrac{1}{2}mv_0^2$ til null:
\[
    \Delta E_\text{K} = 0 - \frac{1}{2}\cdot\SI{20}{\kilo\gram}\cdot(\SI{3.0}{\metre\per\second})^2 = \doubleunderline{\SI{-90}{\joule}}
\]

\smallskip\noindent
\textbf{b)} Det eneste arbeidet på kassen gjøres av friksjonen, $W = -R\cdot s$ (tyngdekraften og normalkraften står vinkelrett på bevegelsen). Arbeid-energi-setningen, $W_\text{total} = \Delta E_\text{K}$, gir
\[
    -R\cdot s = \Delta E_\text{K}
    \quad\implies\quad
    s = \frac{-\Delta E_\text{K}}{R} = \frac{\SI{90}{\joule}}{\SI{40}{\newton}} = \SI{2.25}{\metre} \approx \doubleunderline{\SI{2.3}{\metre}}
\]

\smallskip\noindent
\textbf{c)} Med dobbel startfart blir den kinetiske energien fire ganger så stor ($\SI{360}{\joule}$). Friksjonskraften er den samme, så strekningen blir også fire ganger så lang:
\[
    s = \frac{\SI{360}{\joule}}{\SI{40}{\newton}} = \doubleunderline{\SI{9.0}{\metre}}
\]
}

\oppgaver{Potensiell energi og nullnivå}{k5oppg5}{%
%
En bok med masse $m = \SI{1.2}{\kilo\gram}$ ligger på en hylle \SI{1.8}{\metre} over gulvet. Rett under hylla står et bord som er \SI{0.75}{\metre} høyt.
\begin{enumerate}[a)]
    \item Finn den potensielle energien til boka når vi velger gulvet som nullnivå.
    \item Finn den potensielle energien til boka når vi velger bordplata som nullnivå.
    \item Svarene i a) og b) er forskjellige. Er det et problem? Forklar hva det er som faktisk betyr noe når vi regner med potensiell energi.
\end{enumerate}
%
}

\svar[title={Løsningsforslag}]{%
\noindent
Vi bruker $E_p = m\mathrm{g}h$, der $h$ er høyden over nullnivået.

\smallskip\noindent
\textbf{a)} Med gulvet som nullnivå er $h = \SI{1.8}{\metre}$:
\[
    E_p = \SI{1.2}{\kilo\gram}\cdot\SI{9.81}{\metre\per\second\squared}\cdot\SI{1.8}{\metre} = \SI{21.2}{\joule} \approx \doubleunderline{\SI{21}{\joule}}
\]

\smallskip\noindent
\textbf{b)} Med bordplata som nullnivå er $h = \SI{1.8}{\metre} - \SI{0.75}{\metre} = \SI{1.05}{\metre}$:
\[
    E_p = \SI{1.2}{\kilo\gram}\cdot\SI{9.81}{\metre\per\second\squared}\cdot\SI{1.05}{\metre} = \SI{12.4}{\joule} \approx \doubleunderline{\SI{12}{\joule}}
\]

\smallskip\noindent
\textbf{c)} Det er ikke noe problem. Nullnivået velger vi selv, så verdien av $E_p$ alene har ingen fysisk betydning. Det som betyr noe, er \emph{endringen} i potensiell energi når boka flytter seg. Faller boka for eksempel ned på bordet, er endringen $m\mathrm{g}\cdot\SI{1.05}{\metre} = \SI{12.4}{\joule}$ uansett hvilket nullnivå vi har valgt.
}

\oppgaver{Energi i en fjær}{k5oppg6}{%
%
En fjær med fjærkonstant $k = \SI{250}{\newton\per\metre}$ strekkes $x = \SI{20}{\centi\metre}$ fra likevektsstillingen.
\begin{enumerate}[a)]
    \item Hvor mye elastisk potensiell energi er lagret i fjæra?
    \item Hvor mye må fjæra strekkes for å lagre \SI{20}{\joule}?
\end{enumerate}
%
}

\svar[title={Løsningsforslag}]{%
\noindent
\textbf{a)} Elastisk potensiell energi er $E_p = \tfrac{1}{2}kx^2$:
\[
    E_p = \frac{1}{2}\cdot\SI{250}{\newton\per\metre}\cdot(\SI{0.20}{\metre})^2 = \doubleunderline{\SI{5.0}{\joule}}
\]

\smallskip\noindent
\textbf{b)} Vi løser $E_p = \tfrac{1}{2}kx^2$ for $x$:
\[
    x = \sqrt{\frac{2E_p}{k}} = \sqrt{\frac{2\cdot\SI{20}{\joule}}{\SI{250}{\newton\per\metre}}} = \doubleunderline{\SI{0.40}{\metre}}
\]
Fire ganger så mye energi krever bare dobbelt så stor strekning, fordi $E_p$ er proporsjonal med $x^2$.
}

\oppgaver{Syklist ned en bakke}{k5oppg7}{%
%
En syklist og sykkelen har til sammen masse $m = \SI{90}{\kilo\gram}$. Syklisten starter i ro på toppen av en bakke og triller ned uten å tråkke. Høydeforskjellen er $h = \SI{25}{\metre}$, og farten nederst i bakken er $v = \SI{15}{\metre\per\second}$.
\begin{enumerate}[a)]
    \item Finn den potensielle energien på toppen, med bunnen av bakken som nullnivå.
    \item Finn den kinetiske energien nederst i bakken.
    \item Hvor mye mekanisk energi har gått tapt? Hvor har denne energien blitt av?
    \item Hvor mange prosent av den opprinnelige energien gikk tapt?
\end{enumerate}
%
}

\svar[title={Løsningsforslag}]{%
\noindent
\textbf{a)}
\[
    E_p = m\mathrm{g}h = \SI{90}{\kilo\gram}\cdot\SI{9.81}{\metre\per\second\squared}\cdot\SI{25}{\metre} = \SI{22.1}{\kilo\joule} \approx \doubleunderline{\SI{22}{\kilo\joule}}
\]

\smallskip\noindent
\textbf{b)}
\[
    E_\text{K} = \frac{1}{2}mv^2 = \frac{1}{2}\cdot\SI{90}{\kilo\gram}\cdot(\SI{15}{\metre\per\second})^2 = \SI{10.1}{\kilo\joule} \approx \doubleunderline{\SI{10}{\kilo\joule}}
\]

\smallskip\noindent
\textbf{c)} På toppen er all mekanisk energi potensiell, og i bunnen er all mekanisk energi kinetisk. Tapet er differansen:
\[
    \Delta E = \SI{22.1}{\kilo\joule} - \SI{10.1}{\kilo\joule} = \SI{12.0}{\kilo\joule} \approx \doubleunderline{\SI{12}{\kilo\joule}}
\]
Energien er ikke borte. Den har gått til å varme opp lufta og dekkene, på grunn av luftmotstand og rullemotstand.

\smallskip\noindent
\textbf{d)}
\[
    \frac{\SI{12.0}{\kilo\joule}}{\SI{22.1}{\kilo\joule}} = 0.54 = \doubleunderline{\SI{54}{\percent}}
\]
}

\oppgaver{Friksjon gir varme}{k5oppg8}{%
%
En kloss med masse $m = \SI{2.0}{\kilo\gram}$ glir $s = \SI{3.0}{\metre}$ bortover et horisontalt bord. Den kinetiske friksjonskoeffisienten er $\mu_k = 0.25$.
\begin{enumerate}[a)]
    \item Finn friksjonskraften på klossen.
    \item Hvor mye mekanisk energi blir omdannet til varme på denne strekningen?
    \item Klossen hadde farten $v_0 = \SI{4.0}{\metre\per\second}$ i starten. Hvilken fart har den etter \SI{3.0}{\metre}?
\end{enumerate}
%
}

\svar[title={Løsningsforslag}]{%
\noindent
\textbf{a)} Bordet er horisontalt, så $N = m\mathrm{g}$, og friksjonskraften blir
\[
    R = \mu_k N = 0.25\cdot\SI{2.0}{\kilo\gram}\cdot\SI{9.81}{\metre\per\second\squared} = \doubleunderline{\SI{4.9}{\newton}}
\]

\smallskip\noindent
\textbf{b)} Energien som blir omdannet til varme, er like stor som arbeidet friksjonen gjør mot bevegelsen:
\[
    Q = R\cdot s = \SI{4.9}{\newton}\cdot\SI{3.0}{\metre} = \SI{14.7}{\joule} \approx \doubleunderline{\SI{15}{\joule}}
\]

\smallskip\noindent
\textbf{c)} Klossen starter med
\[
    E_{\text{K},0} = \frac{1}{2}\cdot\SI{2.0}{\kilo\gram}\cdot(\SI{4.0}{\metre\per\second})^2 = \SI{16.0}{\joule}
\]
Etter \SI{3.0}{\metre} er $\SI{14.7}{\joule}$ blitt til varme, så $E_\text{K} = \SI{16.0}{\joule} - \SI{14.7}{\joule} = \SI{1.3}{\joule}$. Farten er da
\[
    v = \sqrt{\frac{2E_\text{K}}{m}} = \sqrt{\frac{2\cdot\SI{1.3}{\joule}}{\SI{2.0}{\kilo\gram}}} = \doubleunderline{\SI{1.1}{\metre\per\second}}
\]
}

\oppgaver{Stein kastet rett opp}{k5oppg9}{%
%
En stein kastes rett opp med startfart $v_0 = \SI{12}{\metre\per\second}$. Vi ser bort fra luftmotstand.
\begin{enumerate}[a)]
    \item Bruk bevaring av mekanisk energi til å finne hvor høyt over utkastpunktet steinen kommer.
    \item Hvor stor fart har steinen når den er \SI{5.0}{\metre} over utkastpunktet?
    \item Du fikk ikke oppgitt massen til steinen. Forklar hvorfor du ikke trengte den.
\end{enumerate}
%
}

\svar[title={Løsningsforslag}]{%
\noindent
Vi velger utkastpunktet som nullnivå. Uten luftmotstand er mekanisk energi bevart.

\smallskip\noindent
\textbf{a)} På toppen står steinen stille, så all kinetisk energi er blitt potensiell energi:
\[
    \frac{1}{2}mv_0^2 = m\mathrm{g}h
    \quad\implies\quad
    h = \frac{v_0^2}{2\mathrm{g}} = \frac{(\SI{12}{\metre\per\second})^2}{2\cdot\SI{9.81}{\metre\per\second\squared}} = \doubleunderline{\SI{7.3}{\metre}}
\]

\smallskip\noindent
\textbf{b)}
\[
    \frac{1}{2}mv_0^2 = \frac{1}{2}mv^2 + m\mathrm{g}h
    \quad\implies\quad
    v = \sqrt{v_0^2 - 2\mathrm{g}h}
\]
\[
    v = \sqrt{(\SI{12}{\metre\per\second})^2 - 2\cdot\SI{9.81}{\metre\per\second\squared}\cdot\SI{5.0}{\metre}} = \doubleunderline{\SI{6.8}{\metre\per\second}}
\]

\smallskip\noindent
\textbf{c)} Massen $m$ står i alle leddene i energilikningen, så vi kan dele den bort. Svaret blir det samme uansett hvor tung steinen er (så lenge vi ser bort fra luftmotstand).
}

\oppgaver{Huske}{k5oppg10}{%
%
Et barn sitter i en huske. Husken trekkes bakover og slippes fra ro når setet er $h = \SI{1.2}{\metre}$ over det laveste punktet. Vi ser bort fra luftmotstand og friksjon.
\begin{enumerate}[a)]
    \item Hvor stor fart har barnet i det laveste punktet?
    \item Hvor høyt over det laveste punktet er barnet når farten er \SI{3.0}{\metre\per\second}?
    \item Snora i husken drar hele tiden i barnet, men gjør ikke noe arbeid. Hvorfor ikke? (Hint: hvilken vinkel er det mellom snorkraften og bevegelsesretningen?)
\end{enumerate}
%
}

\svar[title={Løsningsforslag}]{%
\noindent
Vi velger det laveste punktet som nullnivå.

\smallskip\noindent
\textbf{a)} All potensiell energi blir til kinetisk energi i det laveste punktet:
\[
    m\mathrm{g}h = \frac{1}{2}mv^2
    \quad\implies\quad
    v = \sqrt{2\mathrm{g}h} = \sqrt{2\cdot\SI{9.81}{\metre\per\second\squared}\cdot\SI{1.2}{\metre}} = \doubleunderline{\SI{4.9}{\metre\per\second}}
\]

\smallskip\noindent
\textbf{b)} Energibevaring mellom startpunktet og punktet der farten er $v = \SI{3.0}{\metre\per\second}$:
\[
    m\mathrm{g}h_0 = \frac{1}{2}mv^2 + m\mathrm{g}h
    \quad\implies\quad
    h = h_0 - \frac{v^2}{2\mathrm{g}} = \SI{1.2}{\metre} - \frac{(\SI{3.0}{\metre\per\second})^2}{2\cdot\SI{9.81}{\metre\per\second\squared}} = \doubleunderline{\SI{0.74}{\metre}}
\]

\smallskip\noindent
\textbf{c)} Barnet beveger seg langs en sirkelbue, så bevegelsesretningen er hele tiden langs sirkelen. Snora drar inn mot opphengspunktet, altså vinkelrett på bevegelsen. Da er $\cos 90^\circ = 0$, og snorkraften gjør ikke noe arbeid. Derfor kan vi bruke bevaring av mekanisk energi selv om snora drar i barnet.
}

\oppgaver{Løpe opp trappene}{k5oppg11}{%
%
Eva har masse $m = \SI{70}{\kilo\gram}$ og løper opp trappene fra første til fjerde etasje. Høydeforskjellen er $h = \SI{9.0}{\metre}$, og det tar henne $t = \SI{12}{\second}$.
\begin{enumerate}[a)]
    \item Hvor stort arbeid gjør Eva mot tyngdekraften?
    \item Hvor stor er den gjennomsnittlige effekten hennes?
    \item En vanlig vannkoker har effekt på omtrent \SI{2000}{\watt}. Sammenlikn med svaret i b).
\end{enumerate}
%
}

\svar[title={Løsningsforslag}]{%
\noindent
\textbf{a)} Arbeidet mot tyngdekraften er økningen i potensiell energi:
\[
    W = m\mathrm{g}h = \SI{70}{\kilo\gram}\cdot\SI{9.81}{\metre\per\second\squared}\cdot\SI{9.0}{\metre} = \doubleunderline{\SI{6.2}{\kilo\joule}}
\]

\smallskip\noindent
\textbf{b)}
\[
    P = \frac{W}{t} = \frac{\SI{6180}{\joule}}{\SI{12}{\second}} = \SI{515}{\watt} \approx \doubleunderline{\SI{5.2e2}{\watt}}
\]

\smallskip\noindent
\textbf{c)} Effekten hennes er omtrent en firedel av effekten til en vannkoker. Et menneske kan levere en slik effekt bare i korte perioder.
}

\oppgaver{Bil i konstant fart}{k5oppg12}{%
%
En bil kjører med konstant fart $v = \SI{90}{\kilo\metre\per\hour}$ på flat vei. Luftmotstand og rullemotstand gir til sammen en kraft på $F = \SI{600}{\newton}$ mot bevegelsen.
\begin{enumerate}[a)]
    \item Hvor stor kraft må hjulene skyve bilen fremover med for at farten skal være konstant? Forklar med Newtons 1. lov.
    \item Hvor stor effekt må motoren levere til hjulene?
    \item Hvor mye energi går med til å holde farten konstant i én time?
\end{enumerate}
%
}

\svar[title={Løsningsforslag}]{%
\noindent
\textbf{a)} Farten er konstant, så akselerasjonen er null. Newtons 1.\ lov sier at kraftsummen da må være null. Kraften fremover må derfor være like stor som motstanden: \doubleunderline{$F = \SI{600}{\newton}$}.

\smallskip\noindent
\textbf{b)} Vi bruker $P = F\cdot v$ med $v = \SI{90}{\kilo\metre\per\hour} = \SI{25}{\metre\per\second}$:
\[
    P = \SI{600}{\newton}\cdot\SI{25}{\metre\per\second} = \doubleunderline{\SI{15}{\kilo\watt}}
\]

\smallskip\noindent
\textbf{c)}
\[
    E = P\cdot t = \SI{15000}{\watt}\cdot\SI{3600}{\second} = \doubleunderline{\SI{54}{\mega\joule}}
\]
}

\oppgaver{Sjekk av elastisk støt}{k5oppg13}{%
%
Vogn 1 med masse $m_1 = \SI{1.0}{\kilo\gram}$ ruller med fart \SI{4.0}{\metre\per\second} mot vogn 2 med masse $m_2 = \SI{3.0}{\kilo\gram}$, som står i ro. Etter støtet ruller vogn 1 \emph{tilbake} med fart \SI{2.0}{\metre\per\second}, mens vogn 2 ruller fremover med fart \SI{2.0}{\metre\per\second}.
\begin{enumerate}[a)]
    \item Vis at bevegelsesmengden er bevart. (Husk at hastigheten er negativ når vogna ruller tilbake.)
    \item Vis at den kinetiske energien også er bevart. Hva slags støt er dette?
\end{enumerate}
%
}

\svar[title={Løsningsforslag}]{%
\noindent
Vi velger positiv retning i fartsretningen til vogn 1 før støtet. Da er $v_1 = \SI{4.0}{\metre\per\second}$, $v_2 = 0$, $u_1 = \SI{-2.0}{\metre\per\second}$ og $u_2 = \SI{2.0}{\metre\per\second}$.

\smallskip\noindent
\textbf{a)} Bevegelsesmengde før og etter:
\begin{align*}
    p_\text{før} &= m_1v_1 + m_2v_2 = \SI{1.0}{\kilo\gram}\cdot\SI{4.0}{\metre\per\second} + 0 = \SI{4.0}{\kilo\gram\metre\per\second}\\
    p_\text{etter} &= m_1u_1 + m_2u_2 = \SI{1.0}{\kilo\gram}\cdot(\SI{-2.0}{\metre\per\second}) + \SI{3.0}{\kilo\gram}\cdot\SI{2.0}{\metre\per\second} = \SI{4.0}{\kilo\gram\metre\per\second}
\end{align*}
\doubleunderline{$p_\text{før} = p_\text{etter}$}, så bevegelsesmengden er bevart.

\smallskip\noindent
\textbf{b)} Kinetisk energi før og etter:
\begin{align*}
    E_\text{K,før} &= \frac{1}{2}\cdot\SI{1.0}{\kilo\gram}\cdot(\SI{4.0}{\metre\per\second})^2 = \SI{8.0}{\joule}\\
    E_\text{K,etter} &= \frac{1}{2}\cdot\SI{1.0}{\kilo\gram}\cdot(\SI{2.0}{\metre\per\second})^2 + \frac{1}{2}\cdot\SI{3.0}{\kilo\gram}\cdot(\SI{2.0}{\metre\per\second})^2 = \SI{2.0}{\joule} + \SI{6.0}{\joule} = \SI{8.0}{\joule}
\end{align*}
Både bevegelsesmengde og kinetisk energi er bevart, så støtet er \doubleunderline{elastisk}.
}

\oppgaver{Energitap i et uelastisk støt}{k5oppg14}{%
%
En vogn med masse \SI{2.0}{\kilo\gram} ruller med fart \SI{3.0}{\metre\per\second} og kolliderer med en vogn med masse \SI{1.0}{\kilo\gram} som står i ro. Vognene henger fast i hverandre etter støtet.
\begin{enumerate}[a)]
    \item Finn den felles farten til vognene etter støtet.
    \item Finn den totale kinetiske energien før og etter støtet.
    \item Hvor mye kinetisk energi gikk tapt, og hva ble den omdannet til?
\end{enumerate}
%
}

\svar[title={Løsningsforslag}]{%
\noindent
\textbf{a)} Vognene henger sammen etter støtet og får samme fart $u$. Bevaring av bevegelsesmengde:
\[
    m_1v_1 = (m_1 + m_2)\,u
    \quad\implies\quad
    u = \frac{\SI{2.0}{\kilo\gram}\cdot\SI{3.0}{\metre\per\second}}{\SI{3.0}{\kilo\gram}} = \doubleunderline{\SI{2.0}{\metre\per\second}}
\]

\smallskip\noindent
\textbf{b)}
\begin{align*}
    E_\text{K,før} &= \frac{1}{2}\cdot\SI{2.0}{\kilo\gram}\cdot(\SI{3.0}{\metre\per\second})^2 = \doubleunderline{\SI{9.0}{\joule}}\\
    E_\text{K,etter} &= \frac{1}{2}\cdot\SI{3.0}{\kilo\gram}\cdot(\SI{2.0}{\metre\per\second})^2 = \doubleunderline{\SI{6.0}{\joule}}
\end{align*}

\smallskip\noindent
\textbf{c)} $\SI{9.0}{\joule} - \SI{6.0}{\joule} = \doubleunderline{\SI{3.0}{\joule}}$ kinetisk energi gikk tapt. Den er omdannet til varme, lyd og deformasjon av koblingen mellom vognene. Bevegelsesmengden er bevart, men den kinetiske energien er det ikke. Det er typisk for et uelastisk støt.
}

\oppgaver{Arbeid eller ikke?}{k5oppg15}{%
%
Avgjør om kraften i hvert av tilfellene under gjør et positivt, negativt eller null arbeid. Begrunn svarene.
\begin{enumerate}[a)]
    \item Du holder en tung sekk stille over hodet. (Kraften fra armene dine på sekken.)
    \item Du bærer sekken bortover et flatt gulv med konstant fart. (Kraften fra armene dine, som peker rett opp.)
    \item Du løfter sekken opp fra gulvet. (Kraften fra armene dine.)
    \item En bil bremser. (Friksjonskraften fra veien på bilen.)
    \item En satellitt går i sirkelbane rundt jorda. (Tyngdekraften på satellitten.)
\end{enumerate}
%
}

\svar[title={Løsningsforslag}]{%
\noindent
Arbeid krever både en kraft og en forflytning med en komponent langs kraften: $W = F\cdot s\cdot\cos\theta$.

\smallskip\noindent
\textbf{a)} \doubleunderline{Null.} Sekken flytter seg ikke, så $s = 0$. (Musklene dine blir slitne, men det skyldes prosesser inne i kroppen, ikke arbeid på sekken.)

\smallskip\noindent
\textbf{b)} \doubleunderline{Null.} Kraften peker rett opp, mens forflytningen er horisontal. Vinkelen er $90^\circ$, så $\cos\theta = 0$.

\smallskip\noindent
\textbf{c)} \doubleunderline{Positivt.} Kraften og forflytningen peker begge oppover ($\theta = 0^\circ$). Arbeidet øker den potensielle energien til sekken.

\smallskip\noindent
\textbf{d)} \doubleunderline{Negativt.} Friksjonskraften peker bakover, mens bilen beveger seg fremover ($\theta = 180^\circ$). Arbeidet reduserer den kinetiske energien til bilen.

\smallskip\noindent
\textbf{e)} \doubleunderline{Null.} I en sirkelbane peker tyngdekraften inn mot jordas sentrum, mens bevegelsen hele tiden går langs sirkelen. Kraften står vinkelrett på bevegelsen, så farten til satellitten endrer seg ikke.
}

\oppgaver{Fjærkanon}{k5oppg16}{%
%
En leketøyskanon har en fjær med fjærkonstant $k = \SI{500}{\newton\per\metre}$. Fjæra presses sammen $x = \SI{8.0}{\centi\metre}$, og kanonen skyter en kule med masse $m = \SI{50}{\gram}$ rett opp. Se bort fra luftmotstand og fra den lille høydeendringen mens fjæra strekker seg ut.
\begin{enumerate}[a)]
    \item Hvor mye energi er lagret i fjæra før skuddet?
    \item Hvor høyt over utskytingspunktet kommer kula?
    \item Hvor høyt kommer kula dersom fjæra bare presses sammen halvparten så mye?
\end{enumerate}
%
}

\svar[title={Løsningsforslag}]{%
\noindent
\textbf{a)}
\[
    E_p = \frac{1}{2}kx^2 = \frac{1}{2}\cdot\SI{500}{\newton\per\metre}\cdot(\SI{0.080}{\metre})^2 = \doubleunderline{\SI{1.6}{\joule}}
\]

\smallskip\noindent
\textbf{b)} All elastisk potensiell energi blir til gravitasjonspotensiell energi i toppunktet:
\[
    \frac{1}{2}kx^2 = m\mathrm{g}h
    \quad\implies\quad
    h = \frac{\SI{1.6}{\joule}}{\SI{0.050}{\kilo\gram}\cdot\SI{9.81}{\metre\per\second\squared}} = \doubleunderline{\SI{3.3}{\metre}}
\]

\smallskip\noindent
\textbf{c)} Halvparten så stor kompresjon gir $(\tfrac{1}{2})^2 = \tfrac{1}{4}$ av energien, og dermed en firedel av høyden:
\[
    h = \frac{\SI{3.26}{\metre}}{4} = \doubleunderline{\SI{0.82}{\metre}}
\]
}

\oppgaver{Aking med friksjon}{k5oppg17}{%
%
Mia og akebrettet har til sammen masse $m = \SI{30}{\kilo\gram}$. Hun starter i ro på toppen av en bakke som er $h = \SI{8.0}{\metre}$ høy. Bakken er \SI{40}{\metre} lang, målt langs bakken. I bunnen har hun farten $v = \SI{9.0}{\metre\per\second}$.
\begin{enumerate}[a)]
    \item Hvor stor fart ville hun hatt i bunnen uten friksjon og luftmotstand?
    \item Hvor mye mekanisk energi gikk tapt på vei ned?
    \item Anta at hele energitapet skyldes en konstant friksjonskraft langs bakken. Hvor stor er denne kraften?
\end{enumerate}
%
}

\svar[title={Løsningsforslag}]{%
\noindent
\textbf{a)} Uten tap blir all potensiell energi til kinetisk energi:
\[
    v = \sqrt{2\mathrm{g}h} = \sqrt{2\cdot\SI{9.81}{\metre\per\second\squared}\cdot\SI{8.0}{\metre}} = \SI{12.5}{\metre\per\second} \approx \doubleunderline{\SI{13}{\metre\per\second}}
\]

\smallskip\noindent
\textbf{b)} Energitapet er forskjellen mellom potensiell energi på toppen og kinetisk energi i bunnen:
\[
    \Delta E = m\mathrm{g}h - \frac{1}{2}mv^2 = \SI{2354}{\joule} - \SI{1215}{\joule} = \SI{1.14}{\kilo\joule} \approx \doubleunderline{\SI{1.1}{\kilo\joule}}
\]

\smallskip\noindent
\textbf{c)} En konstant friksjonskraft $R$ langs hele bakken gjør arbeidet $-R\cdot s$. Dette arbeidet er like stort som energitapet:
\[
    R = \frac{\Delta E}{s} = \frac{\SI{1139}{\joule}}{\SI{40}{\metre}} = \doubleunderline{\SI{28}{\newton}}
\]
}

\oppgaver{Heis}{k5oppg18}{%
%
En heis med passasjerer har til sammen masse $m = \SI{800}{\kilo\gram}$. Heisen løftes \SI{30}{\metre} med konstant fart på \SI{20}{\second}.
\begin{enumerate}[a)]
    \item Hvor stort arbeid gjør heismotoren?
    \item Hvor stor effekt leverer motoren?
    \item Regn ut effekten på en annen måte, med $P = F\cdot v$, og sjekk at du får samme svar.
\end{enumerate}
%
}

\svar[title={Løsningsforslag}]{%
\noindent
\textbf{a)} Farten er konstant, så motorkraften er lik tyngdekraften, og arbeidet er økningen i potensiell energi:
\[
    W = m\mathrm{g}h = \SI{800}{\kilo\gram}\cdot\SI{9.81}{\metre\per\second\squared}\cdot\SI{30}{\metre} = \SI{235}{\kilo\joule} \approx \doubleunderline{\SI{2.4e2}{\kilo\joule}}
\]

\smallskip\noindent
\textbf{b)}
\[
    P = \frac{W}{t} = \frac{\SI{235440}{\joule}}{\SI{20}{\second}} = \SI{11.8}{\kilo\watt} \approx \doubleunderline{\SI{12}{\kilo\watt}}
\]

\smallskip\noindent
\textbf{c)} Farten er $v = \SI{30}{\metre}/\SI{20}{\second} = \SI{1.5}{\metre\per\second}$, og kraften er $F = m\mathrm{g} = \SI{7848}{\newton}$:
\[
    P = F\cdot v = \SI{7848}{\newton}\cdot\SI{1.5}{\metre\per\second} = \SI{11.8}{\kilo\watt} \approx \doubleunderline{\SI{12}{\kilo\watt}}
\]
Vi får samme svar som i b).
}

\oppgaver{Ballistisk pendel}{k5oppg19}{%
%
En ballistisk pendel brukes til å måle farten til en kule. En kule med masse $m = \SI{10}{\gram}$ skytes inn i en trekloss med masse $M = \SI{2.0}{\kilo\gram}$ som henger i en snor. Kula blir sittende fast i klossen, og klossen svinger ut slik at den stiger $h = \SI{5.0}{\centi\metre}$.
\begin{enumerate}[a)]
    \item Bruk bevaring av mekanisk energi til å finne farten til klossen (med kula) rett etter støtet.
    \item Bruk bevaring av bevegelsesmengde til å finne farten til kula før den traff klossen.
    \item Hvorfor kan vi ikke bruke bevaring av mekanisk energi for selve støtet?
\end{enumerate}
%
}

\svar[title={Løsningsforslag}]{%
\noindent
Vi kaller farten til kula før støtet $v$, og den felles farten til klossen og kula rett etter støtet $u$.

\smallskip\noindent
\textbf{a)} Etter støtet svinger klossen (med kula) opp. Nå virker bare tyngdekraften og snorkraften, og snorkraften gjør ikke arbeid. Mekanisk energi er derfor bevart:
\[
    \frac{1}{2}(m+M)u^2 = (m+M)\mathrm{g}h
    \quad\implies\quad
    u = \sqrt{2\mathrm{g}h}
\]
\[
    u = \sqrt{2\cdot\SI{9.81}{\metre\per\second\squared}\cdot\SI{0.050}{\metre}} = \doubleunderline{\SI{0.99}{\metre\per\second}}
\]

\smallskip\noindent
\textbf{b)} Under selve støtet er bevegelsesmengden bevart:
\[
    mv = (m+M)\,u
    \quad\implies\quad
    v = \frac{(m+M)\,u}{m}
\]
\[
    v = \frac{\SI{2.01}{\kilo\gram}\cdot\SI{0.990}{\metre\per\second}}{\SI{0.010}{\kilo\gram}} = \SI{199}{\metre\per\second} \approx \doubleunderline{\SI{2.0e2}{\metre\per\second}}
\]

\smallskip\noindent
\textbf{c)} Støtet er uelastisk: kula blir sittende fast i klossen. Mesteparten av den kinetiske energien til kula blir omdannet til varme og deformasjon av treverket. (Her er $E_\text{K}$ før støtet omtrent $\SI{198}{\joule}$, mens den bare er omtrent $\SI{0.99}{\joule}$ rett etter.) Derfor kan vi bare bruke bevaring av bevegelsesmengde for selve støtet, og bevaring av mekanisk energi for svingningen etterpå.
}
